\element{1.4.1}{
  \begin{questionmult}{echantillonnageSpectral}\bareme{haut=2,p=0}
    La fréquence d'échantillonnage $f_e$, la durée de l'acquisition $T$, la résolution spectrale $\Delta f$ et l'étendue spectrale $f_\text{max}$ sont reliées par
    \begin{multicols}{2}
      \begin{reponses}
        \bonne{$f_\text{max}=\frac{f_e}{2}$}
        \bonne{$\Delta f=\frac{1}{T}$}
        \mauvaise{$\Delta f=\frac{f_e}{2}$}
        \mauvaise{$f_\text{max}=\frac{1}{T}$}
      \end{reponses}
    \end{multicols}
  \end{questionmult}
}

\element{1.4.2}{
  \begin{question}{shannon}\bareme{1}
    Le théorème de Shannon dit que
    \begin{multicols}{2}
      \begin{reponses}
        \bonne{Il faut échantillonner à une fréquence d'au moins $2f_{max}$}
        \mauvaise{Il faut échantillonner à une fréquence d'au plus $2f_{max}$}
        \mauvaise{Il faut échantillonner à une fréquence d'au moins $\frac{f_{max}}{2}$}
        \mauvaise{Il faut échantillonner à une fréquence d'au plus $\frac{f_{max}}{2}$}
      \end{reponses}
    \end{multicols}
  \end{question}
}
